\section{Security Games and Reduction Pseudocode}
\label{app:security-games}

This appendix gives the full security experiments used in the paper and
expands the reduction from \Cref{sec:verified-reduction} into
executable-style pseudocode.

\subsection{The IND-CPA experiment}
\label{app:ind-cpa-game}

IND-CPA is the standard confidentiality notion for public-key encryption.
We use its multi-query left-or-right formulation: the challenger samples one
hidden bit \(b\), and every encryption query selects its plaintext with that
same bit.
The evaluation key is public, so the adversary can evaluate ciphertexts
locally without a challenger oracle.

\begin{algorithm}
\caption{Multi-query left-or-right IND-CPA experiment}
\label{alg:ind-cpa-game}
\begin{algorithmic}
  \Procedure{Game}{$\cB$}
    \State $b \gets \{0,1\}$
    \State $(pk,evk,sk) \gets KeyGen()$
    \State $b^* \gets \cB^{OEnc_{\mathsf{CPA}}}(pk,evk)$
    \State \Return $(b^*=b)$
  \EndProcedure
  \State
  \Procedure{$OEnc_{\mathsf{CPA}}$}{$m_0,m_1$}
    \State $c \gets Enc(pk,m_b)$
    \State \Return $c$
  \EndProcedure
\end{algorithmic}
\end{algorithm}

The adversary may choose several encryption queries adaptively.
We state security in winning-probability form: the scheme's IND-CPA assumption
supplies a bound
\[
 \WinProb[\mathsf{IND\text{-}CPA}_{S}(\cB)]
 \leq \beta_{\mathsf{CPA}}(\cB)
\]
for every adversary \(\cB\) with the required oracle interface.
Equivalently, if advantage means excess winning probability over random
guessing, then
\(\beta_{\mathsf{CPA}}(\cB)=\tfrac12+\mathsf{Adv}_{\mathsf{CPA}}(\cB)\).

\subsection{The IND-CPAD experiment}
\label{app:ind-cpad-game}

Li and Micciancio introduced IND-CPA with decryption, or IND-CPAD, to model
attacks that exploit released approximate decryptions~\cite{lm}.
The game retains the hidden bit, encryption oracle, and final guess of
IND-CPA, and adds homomorphic evaluation and restricted decryption.
Its challenger maintains a table
\[
 T=[(m_0,m_1,c),\ldots].
\]
Each row records the actual ciphertext \(c\) together with the plaintext it
would represent in the left challenge world and in the right challenge world.
Only the ciphertext for the world selected by \(b\) is sampled.

\begin{algorithm}
\caption{IND-CPAD Challenger Skeleton}
\label{alg:ind-cpad-game}
Ingredients:
\begin{itemize}
  \item An approximate homomorphic encryption scheme
    $S=(KeyGen,Enc,Eval,Dec)$.
  \item A decrypt-query bound $q\in\bbN$.
\end{itemize}
Global challenger state:
\begin{itemize}
  \item keys $pk,evk,sk$;
  \item hidden challenge bit $b\in\{0,1\}$;
  \item table $T\in(M\times M\times C)^*$ of challenge plaintext pairs and
    produced ciphertexts;
  \item decrypt counter $d$.
\end{itemize}
\begin{algorithmic}
  \Procedure{Game}{$\cA$}
    \State $b \gets \{0,1\}$
    \State $(pk,evk,sk) \gets KeyGen()$
    \State $T \gets [\,]$
    \State $d \gets 0$
    \State $b^* \gets \cA^{OEnc,OEval,ODec}(pk,evk)$
    \State \Return $(b^*=b)$
  \EndProcedure
\end{algorithmic}
\end{algorithm}

The encryption oracle appends a new row.
The evaluation oracle applies the same gate separately to the recorded
left-world and right-world plaintexts, then stores both results alongside the
evaluated ciphertext.
The decryption oracle returns a plaintext only if the two recorded plaintexts
agree.
This restriction is what makes decryption compatible with a left-or-right
game: if the two plaintexts differed, releasing the decryption could reveal
the hidden bit immediately.

\begin{algorithm}
\caption{IND-CPAD Oracles}
\label{alg:ind-cpad-oracles}
The oracles share the challenger state from \Cref{alg:ind-cpad-game}.
\begin{algorithmic}
  \Procedure{OEnc}{$m_0,m_1$}
    \State $c \gets Enc(pk,m_b)$
    \State $T \gets T \doubleplus [(m_0,m_1,c)]$
    \State \Return $c$
  \EndProcedure
  \State
  \Procedure{OEval}{$f,\mathcal{I}$}
    \State $\vec r \gets (T_i)_{i\in\mathcal{I}}$
    \State $m_0' \gets f((r.m_0)_{r\in\vec r})$
    \State $m_1' \gets f((r.m_1)_{r\in\vec r})$
    \State $c' \gets Eval(f,(r.c)_{r\in\vec r})$
    \State $T \gets T \doubleplus [(m_0',m_1',c')]$
    \State \Return $c'$
  \EndProcedure
  \State
  \Procedure{ODec}{$i$}
    \State \textbf{assert} $d<q$
    \State $d \gets d+1$
    \State \textbf{assert} $i<|T|$
    \If{$T_i.m_0\ne T_i.m_1$}
      \State \Return $\mathsf{None}$
    \Else
      \State $y\gets\mathsf{SampleCurrent}(T_i)$
      \State \Return $\mathsf{Some}(y)$
    \EndIf
  \EndProcedure
\end{algorithmic}
\end{algorithm}

The adversary may interleave encryption, evaluation, and decryption calls
adaptively.
Calls beyond the \(q\)-query decryption budget, out-of-range table indices,
and other failed assertions abort that execution.
In SSProve, an aborted execution contributes no probability to winning.

Here \(\mathsf{SampleCurrent}\) is instantiated by the surrounding game.
For a row whose two recorded plaintexts agree,
\(T_i=(m,m,\mathsf{Some}(\bar c,e))\), the real game uses
\(\mathsf{Flood}_e(\mathsf{dec}_0(sk,T_i.c))\), whereas the
simulated-decryption game uses \(\mathsf{Flood}_e(m)\), exactly as in
\Cref{eq:real-sim-decrypt}.
The formal proof factors the common prefix of \textsc{ODec} from this final
sampling continuation.

\subsection{Reduction adversary}

\begin{algorithm}
\caption{IND-CPA reduction adversary $\mathcal{B}_{\cA,q}$}
\label{alg:ind-cpa-reduction}
This is the conventional reduction endpoint used in the main proof.
It stores the same logical table as the IND-CPAD challenger, but obtains
challenge encryptions from the outer IND-CPA encryption oracle.
\begin{algorithmic}[1]
  \Procedure{$\mathsf{Flood}$}{$m,e$}
    \State $\sigma_e\gets\max\{1,e\}\gamma$
    \State $\eta \gets D_{\bbZ^n,0,\sigma_e^2}$
    \State \Return $J_m(\eta)$
  \EndProcedure
  \Statex
  \Procedure{$\mathcal{B}_{\cA,q}$}{$pk,evk$}
    \State $T \gets [\,]$
    \State $d \gets 0$
    \State $b^* \gets
      \cA^{OEnc_{\mathcal{B}},OEval_{\mathcal{B}},ODec_{\mathcal{B}}}(pk,evk)$
    \State \Return $b^*$
  \EndProcedure
  \Statex
  \Procedure{$OEnc_{\mathcal{B}}$}{$m_0,m_1$}
    \State $c \gets OEnc_{\mathsf{CPA}}(m_0,m_1)$
      \Comment{outer IND-CPA encryption oracle}
    \State $T \gets T \doubleplus [(m_0,m_1,c)]$
    \State \Return $c$
  \EndProcedure
  \Statex
  \Procedure{$OEval_{\mathcal{B}}$}{$f,\mathcal{I}$}
    \State $\vec r \gets (T_i)_{i\in\mathcal{I}}$
    \State $m_0' \gets f((r.m_0)_{r\in\vec r})$
    \State $m_1' \gets f((r.m_1)_{r\in\vec r})$
    \State $c' \gets Eval(f,(r.c)_{r\in\vec r})$
    \State $T \gets T \doubleplus [(m_0',m_1',c')]$
    \State \Return $c'$
  \EndProcedure
  \Statex
  \Procedure{$ODec_{\mathcal{B}}$}{$i$}
    \State \textbf{assert} $d<q$
    \State $d \gets d+1$
    \State \textbf{assert} $i<|T|$
    \State $(m_0,m_1,c) \gets T_i$
    \If{$m_0\ne m_1$}
      \State \Return $\mathsf{None}$
    \EndIf
    \State \textbf{assert} $c=\mathsf{Some}(\_,e)$
    \State $y \gets \mathsf{Flood}(m_0,e)$
    \State \Return $\mathsf{Some}(y)$
  \EndProcedure
\end{algorithmic}
\end{algorithm}

\FloatBarrier

The reduction does not know the IND-CPA challenge bit.
This is harmless for decryption queries that return a value: the simulator
only returns a flooded plaintext when the two recorded plaintexts agree,
$m_0=m_1$, so there is no bit-dependent plaintext choice left to make.
